Imagine you are an administrator at a Turkish university.
A new internal policy lands on your desk. Before it can go into effect,
you need to verify that it complies with the regulations set by YÖK ---
the Higher Education Council that governs all Turkish universities.
That means reading through dozens of pages of dense legal text,
cross-referencing multiple regulation documents, and making judgment calls
about ambiguous wording. This process can take hours, and it has to be
repeated for every new policy.

Now imagine a system that does this automatically.

That is exactly what we built. \textbf{ComplianceAI} is a
retrieval-augmented generation (RAG) system that takes a university policy
document, searches a database of nine official YÖK regulatory documents for
the most relevant passages, and asks a large language model to decide
whether the policy is \textit{compliant}, \textit{partially compliant}, or
\textit{non-compliant} --- all in under two seconds.

\subsection*{Why RAG?}

A natural first question is: why not just ask GPT-4o directly?
The answer is that language models are trained on large amounts of text from
the internet, but they are not guaranteed to know the specific wording of a
particular YÖK regulation from 2024 --- especially the fine-grained numerical
thresholds and procedural rules that determine compliance.
RAG addresses this by retrieving the actual regulation text at query time
and feeding it to the model as context \citep{Sun2025, Hillebrand2024}.
This way, the model reasons about what the regulation \textit{actually says},
not what it \textit{thinks} it says.

To test whether this matters, we designed our evaluation specifically around
cases where the model would be likely to fail without retrieval.
As we show in Section~\ref{sec:results}, the results are clear:
RAG-based configurations outperform the no-RAG baseline by \textbf{10 percentage
points} on our benchmark.

\subsection*{What We Contribute}

This paper makes the following contributions:

\begin{itemize}
  \item We present \textbf{ComplianceAI}, a complete end-to-end system for
    automated regulatory compliance checking in Turkish higher education,
    including a FastAPI backend, a React dashboard, and five retrieval pipeline
    variants (no-RAG, BM25, Dense, Hybrid, and Multi-Query).

  \item We build the \textbf{first annotated benchmark dataset} for this domain:
    40 compliance assessment cases specifically designed to require access to
    YÖK regulation text, with balanced class distribution (13 compliant,
    14 partial, 13 non-compliant).

  \item We conduct a \textbf{systematic comparison} of all five pipeline variants
    with two LLM backends (GPT-4o and GPT-4o-mini), finding that Hybrid and
    BM25 retrieval consistently outperform other approaches, and that
    GPT-4o-mini matches GPT-4o accuracy at 19$\times$ lower cost.

  \item We run \textbf{ablation studies} on chunk size and retrieval depth,
    identifying 350-word chunks and Top-5 retrieval as the optimal settings
    for this task.

  \item We document an important \textbf{failure mode}: all models
    systematically under-predict the non-compliant class, defaulting to
    ``partial'' even for clear violations --- a finding that should
    inform any real-world deployment of similar systems.
\end{itemize}

The rest of this paper is structured as follows.
Section~\ref{sec:related} reviews related work.
Section~\ref{sec:architecture} describes the system.
Section~\ref{sec:methodology} explains how we built the benchmark and ran
the experiments.
Section~\ref{sec:setup} covers implementation details.
Section~\ref{sec:results} presents and discusses the results ---
the most important part, with five tables and four figures.
Section~\ref{sec:conclusion} wraps up with limitations and future directions.
